\subsection{The next obstruction for the recorded degree-seven action}

\begin{statement}{D.32}{Proved}{An exhaustive degree-eight obstruction for the recorded representative}
Let $\alpha_7$ be the specific automorphism of
$T_2=B/D_8(K)$ recorded in D.31, and put
\[
 T_3=B/D_9(K),\qquad
 Z_8=D_8(K)/D_9(K)\simeq C_2^{8230}.                 \tag{D.32.1}
\]
The natural map $T_3\to T_2$ has kernel $Z_8$ and
$|T_3|=576\cdot2^{11563}$. The recorded $\alpha_7$
has no lift to $T_3$ that sends $s,r$ to conjugates
of $s,r$ and sends the cusp $t=sr$ to a power
compatible with its recorded lower image.

This is a complete test for that lower action.
The fixed-conjugator correction map on $Z_8$ has
rank 7753. The multiplier contributes one more
independent column. The full lower centralizer is
\[
 C_{N/D_8(K)}(s)\simeq C_2^{1666};                   \tag{D.32.2}
\]
its degree-four through degree-seven dimensions are
$37,102,348,1179$. Adding all its defects gives rank
7777. An explicit linear functional annihilates that
entire image and evaluates to one on the recorded
cusp error.

Moreover, $\alpha_7$ has no automorphism lift to
the full relative pro-two group $\mathcal P_2$
that maps $t$ to a power of itself. Consequently the specific compatible
quotient action recorded in D.31 has no genuine full
profinite lift. Its quotient automorphisms remain
valid. The attainability of the same CM pair by a
different action is still undecided.
\end{statement}
\begin{proof}
\emph{The exact error and the next layer.}
The seed program reconstructs the recorded D.31
conjugators, including every upper central and
triple-commutator correction in their stated order.
It evaluates the full $4,903,248$-letter cusp word
in the rank-four Magnus algebra through degree
eight. Degrees one through seven vanish. The
degree-eight coefficient has 32,296 nonzero
monomials. No change of the lower representative
enters this evaluation.

The restricted Hall basis in
\cite[Proposition 4.4]{dEfrat2023} consists in this
degree of 8160 ordinary degree-eight Lyndon brackets,
60 squares of degree-four brackets, six fourth
powers of degree-two brackets, and four eighth
powers of the generators. Its dimension is 8230.
The dimension sum through degree eight is 11568;
multiplying by $|Q_0|=18$ proves the order in the
statement. Conjugation on this central $K$-layer
factors through $Q_0$. Since the two recorded
conjugators lie in $K$, their corrections have
precisely the map
\[
 (b,a)\longmapsto(1-s)b+s(1-r)a,\qquad b,a\in Z_8.
                                                            \tag{D.32.3}
\]
Its rank is 7753.

The cusp has order 48 in $T_2$ and order 96 in
$T_3$. Indeed the nonzero degree-one term of $t^6$
has eighth power in degree eight and sixteenth
power in degree sixteen. Direct evaluation of
$t^{48}$ gives 6561 nonzero degree-eight monomials.
Every multiplier change retaining the lower action
therefore contributes either zero or this one
column. The change by 144 is compatible also with
the recorded finite-shadow multiplier and has the
same column, since $144\equiv48\pmod{96}$. Adding
it raises the rank to 7754; the target remains
outside. This includes the multiplier freedom
that first becomes visible in this transition.

\emph{The retained homology quotient.}
In $K/\Phi(N)$ the quotient by $N/\Phi(N)$ is
$O_2(Q)=V_4\times Q_8$, of class two and exponent
four. The image of $\gamma_3K$ in the elementary
abelian module $N/\Phi(N)$ is generated by the
conjugates of the 64 words
$[a_i,[a_j,a_k]]$, with $1\leq i,j,k\leq4$.
Successive multiplication by the four augmentation
operators gives the exact dimensions
\[
 87,\ 67,\ 43,\ 19,\ 4,\ 0
 \quad\text{for }\gamma_jK,\quad3\leq j\leq8.       \tag{D.32.4}
\]
The other factors of $D_8(K)$ are fourth powers
from $\gamma_2K$, squares from $\gamma_4K$, and
eighth powers from $K$. Their half-powers already
lie in $N$, so they vanish in $N/\Phi(N)$. Thus
$D_8(K)\subset\Phi(N)$ and the natural map from
$B$ to $G_2$ factors through $T_2$.

\emph{Exhausting the lower centralizer.}
The code uses the exact involution on the chosen
free basis of $K$:
\[
 a_1\mapsto a_1^{-1},\quad a_2\mapsto a_2^{-1},\quad
 a_3\mapsto a_4^{-1},\quad a_4\mapsto a_3^{-1}.
                                                            \tag{D.32.5}
\]
Its Magnus substitution retains every coefficient
through degree eight. It is checked against the
actual Bass--Serre words.

For $n\geq3$, the kernel of
$N/D_{n+1}(K)\to N/D_n(K)$ is the full restricted
layer $Z_n=D_n(K)/D_{n+1}(K)$. It is central in
the former group. The connecting map
\[
 C_{N/D_n(K)}(s)\longrightarrow Z_n/(1-s)Z_n
                                                            \tag{D.32.6}
\]
which sends an element to its commutation defect
is therefore a genuine group homomorphism. The
program maintains a full polycyclic sequence for
the source, including all its central coordinates.
It eliminates the connecting columns in descending
sequence order. Each remaining generator has its
own earlier leading coordinate and uses only later
pivot generators for its correction. The resulting
products generate the complete kernel: their
filtered leading coordinates are independent, so
their ordered products give at least the required
number of elements, while the kernel order is
exactly the source order divided by the computed
image order. Lifting these products and adjoining
$\ker(1-s:Z_n\to Z_n)$ gives the entire next
centralizer.

The initial centralizer in $N/D_3(K)$ has dimension
six. The successive exact counts are
\[
\begin{array}{c|rrrrr}
 n&3&4&5&6&7\\\hline
 \dim Z_n&20&70&204&690&2340\\
 \dim Z_n^s&11&40&105&355&1179\\
 \operatorname{rank}(\text{connecting map})&1&6&3&13&7\\
 \log_2|C_{N/D_{n+1}(K)}(s)|&16&50&152&494&1666
\end{array}                                                   \tag{D.32.7}
\]
Every lifted generator is directly checked fixed
at its asserted level. The final generators all
have degree at least four. Their products commute
and their squares vanish modulo $D_8(K)$, proving
(D.32.2), including its exponent and the displayed
degree distribution.

Changing the right $s$-conjugator by these elements
adds their degree-eight defects. The preceding
conjugator belongs to $K$, which acts trivially on
$Z_8$, so no further conjugation factor is suppressed.
This exhausts the finite conjugator choices as well.
An $s$-conjugator for any proposed lift can be
normalized into $F$ by multiplying it by $s$, since
its image in $S_3$ centralizes the transposition.
Its lower difference from the recorded conjugator
then belongs to $F$ and centralizes $s$ in $T_2$.
Projecting this difference to $G_2$ and applying
D.25's projected-centralizer calculation puts it
in $N/D_8(K)$. Thus (D.32.2) includes every such
lower difference, not only ones initially chosen
inside $N$.
The defect map into the cokernel of (D.32.3) is a
group homomorphism. All 1666 generators are included.
The order-three centralizer lifts across a finite
two-kernel, by conjugacy of the two order-three
complements, and contributes no additional classes.
The final rank is 7777. The certificate contains a
functional on the 65536-dimensional degree-eight
tensor space, verifies its vanishing on every row
of the image basis and every lower-centralizer
defect, and verifies its value one on the target.
This proves the negative finite conjugator test.
The parametrization also accounts for generation.
Its tuples define homomorphisms from $B$ that induce
the recorded lower action and the identity on $Q_0$.
The relative pro-two universal property extends
them to $\mathcal P_2$, and the characteristic
filtration gives the induced maps of $T_3$.
Their maps on $K/D_2(K)$ are the recorded
isomorphism, so the finite Frattini theorem makes
them surjective, hence automorphisms. The proved
emptiness occurs before any further generation
condition is needed.

\emph{Promotion to the actual relative pro-two source.}
Let $S_0$ be the preimage of $\langle s\rangle$
in $B\to Q_0$. Its index is nine. The Bass--Serre
quotient has nine edges, six vertices over the
order-two factor, and three over the order-three
factor. Three of the first six vertices have
stabilizer $C_2$, and all other stabilizers are
trivial. Its first Betti number is one. Thus
\[
 S_0\simeq C_2*C_2*C_2*\mathbb Z,\qquad
 S\simeq C_2\amalg_2 C_2\amalg_2 C_2\amalg_2\mathbb Z_2,
                                                            \tag{D.32.8}
\]
where $S$ is its closure in $\mathcal P_2$.
The second map is exactly the pro-two completion:
$K$ has index two in $S_0$, and the normal core
in $B$ of any normal subgroup of $K$ with finite
two-power index still has two-power index in $K$.
Consequently the
relative topology induces the full pro-two topology
on $S_0$.

The original \cite[\S2, Theorem A$'$ and Theorem
B$'$]{A_HerfortRibes1985} applies with the class of
finite two-groups. It proves that the involutions
of $S$ are conjugate into the three vertex factors,
whose generators are themselves discrete
$B$-conjugates of $s$. As $S$ is a Sylow pro-two
subgroup, every involution of $\mathcal P_2$ is
therefore conjugate to $s$. The same graph also
gives $C_{\mathcal P_2}(s)=\langle s\rangle$:
the normalizer quotient of $S$ is $C_3$, acting
cyclically on its three distinct involution classes;
an element centralizing $s$ projects into
$C_{Q_0}(s)$ and hence normalizes the full preimage
$S$ of $\langle s\rangle$. Its action must fix the
class of $s$, so its image in this $C_3$ is trivial.
It therefore lies in $S$, where
Theorem B$'$ computes its centralizer.

Now an automorphism lifting $\alpha_7$ must send
$s$ to a conjugate. Its image of $r$ is conjugate
to $r$ with the prescribed lower conjugator,
because the kernel of $\mathcal P_2\to T_2$ is
pro-two and the order-three complements are
conjugate. Normalize the $s$-conjugator into the
preimage $F$ of the identity in $S_3$ by multiplying
it by $s$ when necessary. The corresponding
conjugator of D.31 is already in $F$. Their lower
difference centralizes $s$. Because $T_2$ retains
$G_2$, D.25's projected-centralizer calculation
forces this difference into $N/D_8(K)$. The
complete test (D.32.2)--(D.32.7) therefore applies
to every such lift, giving a contradiction.
The finite maps and the full source have now been
related explicitly. This excludes the recorded
representative, while leaving changes to its
lower action within the same CM pair available
for further testing.
\end{proof}
